\section*{Introduction}
\dropcap{S}egmenation of cells or their nuclei is an evolving computational field, currently dominated by machine learning as the prevalent tool and foundation of most algorithms. The project behind this report focused on evaluating a selection of segmentation algorithms, including a non-ML model, using four benchmarks. All algorithms will receive a short introduction followed by their evaluation and comparison to each other. Afterwards, a more in-depth presentation of a custom pipeline for preparation, segmentation and evaluation of Xenium data concludes this report.\\
The following colour-codes apply: \inlinepath{path or file}, \inlinecode{code or file-content}, \inlineprogram{program or command}.
%
% Important to consider is that the datasets generated by the workflows are high-data thus trying to analyize anything is computationally expensive. especially identification of single cells.\\
% if any feature is inaccurately assigne d otr not present at all this cwill have consequences further downstreenm. \\
% -----------------------------
% the report below can be conceptualized as investigating a task-specific challenge. the segmentation and evaluation are constraint to fluorescet spleen samples and spatially resolved transcriptomics. no image maniputatiopn is done on the data. the preparations areexplaind in [pipeline]
% 
\subsection*{\textsc{Xenium}}
A device for \textit{in-situ} spatially resolved transcriptomics (SRT) developed by \textsc{10XGenomics}. The multiplex spatial analysis tool allows single cell resolution of mRNA-location with additional morphological stains. The samples are either formalin fixed \& paraffin embedded (FFPE) or fresh frozen \cite{janesick_high_2023}.\\
It has a high multiplex panel with high spatial resolution in XYZ planes --- presumably, using a confocal microscope. The panel is comprised of 479 genes, 379 are integrated, provided by \textsc{10XGenomics}, 100 are customized for immune cells (see Supplement XYZ for a complete list).\\ % [panel: in chat with anika 13th april]
The DAPI stain, as part of the morphological recording, is measured in all planes and available as a multi-layer image, as well as a max-projection in \inlinepath{morphology\us focus\us 0000.ome.tif}. ATP1A1/E-Cadherin/CD45, as the boundary stain, 18S rRNA and alphaSMA/Vimentin, as the interior stains, are collected in the \inlinepath{morphology\us focus} multi-file image, as single plane measurements, as well. \\
The imaging mechanism of the target mRNA is based on rolling circle amplification, extended by a padlock-probe. The padlock-probe, a circular oligo-nucleotide with a complimentary sequence to the target, hybridizes to the mRNA, a polymerase then replicates the sequence of the padlock (RCA). The replicated sequence is targeted by the fluorophores, allowing for a high-intensity recording. If the padlock-probe cannot hybridize correctly, the instability hinders fluorphore hybridisation to the amplified sequence. \cite{jain_padlock_2021}\\
After decoding of the transcript-location, and other post-processing steps, the data-output is as below, listing those important for the segmentation algorithms.
% 
\vspace{9pt}
\dirtree{%
    .1 \textless Xenium output directory\textgreater.
    .2 cell\us boundaries.parquet.
    .2 morphology.ome.tif.
    .2 morphology\us focus.
    .3 morphology\us focus\us 000\{0..3\}.ome.tif.
    .2 nucleus\us boundaries.parquet.
    .2 transcripts.parquet.
}
% 
\subsection*{Spleen}
The Spleen, located behind the diaphragm, is the largest lymphoid, blood-filtering organ in both mice and humans. Apart from recycling iron from old red blood cells, it has important immunological functions. The macrostructures white and red pulp (WP \& RP), separated by the peri-follicular zone (RFZ), serve unique and distinct roles in local and systemic immune-response. Contrary to lymph nodes the spleen is integrated into the immune-system via blood-vessels, as it lacks afferent lymphatic-vessels. \cite{mebius_structure_2005} \cite{cesta_normal_2006} \cite{bronte_spleen_2013} \cite{lewis_structure_2019}\\
The primary function of the RP is extracting and recycling non-functioning red blood cells (RBCs), be it because of age, damage or anti-body markers. Further, the diversity of leukocytes present in the RP enable immune functions like anti-body distribution into the blood-stream.\\
The WP, marginally segregated from the RP by the RFZ, principally constituted of T and B cells, functions as the draining secondary lymphoid organ of the blood. Further segmented in T and B cell-domains, the WP is (place) for B cell isotype-switching and somatic hyper-mutation, in the B cell zones, and location for T cell-B Cell or T cell-dentritic cell interactions, in the T cell zone.\\
The complexity of the spleen-tissue in combination with the, though diverse in type, homogeneous cell-shape present a unique tissue, expecially for computer vision challenges.
\subsection*{Benchmarking}
Computer predictions of shapes is conventionally evaluated by comparing the masks to a human-drawn ground-truth. The quantification of the prediction quality is facilitated through similarity analysis. Precision, recall, accuracy or Jaccard-index, give a numeric value to the similarity between a prediction and a \emph{truth}. These values are representations of \textit{True Positives}, \textit{True Negatives}, \textit{False Positives} and \textit{False Negatives}. Defining these metrics in computer vision is constraint to human-labelled ground-truths, and must be calculate through a separate metric, the Jaccard-index. This, naturally, is not constrained to cell predictions --- MRI and CT predictions or financial data can similarly be assessed using JI and it's derived values \cite{durkee_generalizations_2023}. 
\subsubsection*{Jaccard-Index}
To measure the similarity of two sets, or, less abstract, the overlap of to areas, the intersection and union of these two sets is considered. The Jaccard-index (JI) \eqref{eq:jaccard}, also called Intersection over Union, combines the information on the sets of TP, FP and FN, hence it is used to categorise retrieved samples, cell predictions in this instance, as \textit{true} or \textit{false} depending on a fixed threshold.
\begin{equation}
\label{eq:jaccard}
    J = \frac{SR\cap GT}{SR \cup GT} = \frac{TP}{TP+FP+FN} 
\end{equation}
\subsubsection*{Precision and Recall}
Given a set of samples, one subset categorised as relevant, the other not, the precision \eqref{eq:PR} is the proportion between relevant, retrieved samples and all retrieved samples. The precision decreases with increasing irrelevant, retrieved samples, or decreasing relevant, retrieved samples.
\begin{equation}
\label{eq:PR}
    PR = \frac{TP}{TP+FP}
\end{equation}
Given the set as before, the recall \eqref{eq:RE} is the proportion between relevant, retrieved samples and all relevant samples. The recall decreases with decreasing retrieved, relevant samples.
\begin{equation}
\label{eq:RE}
    RE = \frac{TP}{TP+FN}
\end{equation}
\subsubsection*{F-score}
The F-score, or F-measure, is the arithmetic mean of precision and recall. The first arithmetic mean is abbreviated with F1 and stands as the default value for benchmarking segmentations.
\begin{equation}
\label{eq:F1}
    F1 = \frac{2}{\frac{1}{PR}+\frac{1}{RE}} = 2\frac{PR\cdot RE}{PR+RE} = \frac{2TP}{2TP+FP+FN}
\end{equation}
It's value succinctly represents the quality of a segmentation by combining the information of true and false predictions.
\subsubsection*{Dice-Sørensen Coefficient}
An alternative to calculating a relation between precision and recall, independent of a threshold, contrary to the F-score, is the DCE value \eqref{eq:DCE}. The relation between TP, FP and FN remains the same as for F1. 
\begin{equation}
\label{eq:DCE}
    DCE = \frac{2|SR\cap GT|}{|SR|+|GT|} = \frac{2TP}{2TP+FP+FN}
\end{equation}
The threshold independence suits well for a pixel-level evaluation of segmentation masks, as opposed to the object oriented measure of the F-score.
\subsubsection*{panoptic quality and SEG}
The \inlineprogram{DeepCell-Toolbox} additionally includes the values SEG \eqref{eq:SEG} and panoptic quality (PQ) \eqref{eq:PQ}. The panoptic quality was proposed to measure the quality of the labelling task, the semantic segmentation, and the ability to detect the shape of an object, the instance segmentation \cite{kirillov_segment_2023}. 
\begin{equation}
\label{eq:PQ}
    PQ = \frac{|TP|}{|TP| + \frac{1}{2}|FP| + \frac{1}{2}|FN|} \times \frac{\sum_{(y,\hat y)\epsilon TP}IoU(y, \hat y)}{|TP|}
\end{equation}
The SEG value, developed for cell-tracking challenges, considers only those as matching if the intersection of segmentation result and ground-truth is greater than half the ground-truth area. The final SEG score is the average of all IoU values\cite{maska_benchmark_2014}.
\begin{equation}
\label{eq:SEG}
    SEG = \frac{|SR \cap GT|}{|SR\cup GT|}, when |SR \cap GT| > 0.5\cdot|GT|
\end{equation}
\\
The ground truth is either prepared by a single expert or undergoes multiple iteration of verification of variable expert-composition, e.g. crowd-sourced approach for the \textsc{TissueNet} project. These masks of course exhibit inconsistent quality, repeated analysis has shown a constraint to a 0.5\% inter-annotator-disagreement. \\
Evaluating and training Large Language Models (LLM) might primarily rely on machine-jury/-annotators, as has been suggest in \cite{calderon_alternative_2025}. This is of course not entirely applicable to medical or cell-segmentation. Other reference-free attempts have been made to accommodate the increasing aptness of computer segmentation. The principal component analysis (PCA) of \inlineprogram{CellSegmentationEvaluation} relies on metrics that apply to both the staining-image and the prediction and compares these \cite{chen_evaluation_2023}. Another approach is instead probability based, called \inlineprogram{SEG}, not to be confused with the SEG-value. Using multiple algorithms to predict cell-instances and accumulating their results in a probability map, one can evaluate the applied set of algorithms without a human-annotator \cite{sims_seg_2023}.\\